\documentclass{article}

\usepackage{tikz-cd, amsmath, amssymb, amsthm, mathtools, verbatim}

\newtheorem{lemma}{Lemma}
\newtheorem{conjecture}{Conjecture}
\newtheorem{sketch}{Proof Sketch}

\usepackage[utf8]{inputenc}
\usepackage[T1]{fontenc}
\usepackage{lmodern} % Better standard font

\begin{document}

\section{Polynomial Functors}

Give semantics to W-Types, which should be enough for our case (hopefully). By encoding a set of operations an their arities, which we can think of as the type constructors. Where B are the operations and A the arities, we each also fibre them over a set of types. We only consider endofunctors here.

\begin{center}
    \begin{tikzcd}
        & A \ar[r] \ar[dl] & B \ar[dr] & \\
        I & & & I
    \end{tikzcd}
\end{center}

\section{Given Constructors}

We need the subtype to be stable under these, so we need an $F-Alg$ for the polynomial functor build from the constructors.

\begin{center}
    \begin{tikzcd}
        F(\Sigma_{x:X} P(x)) \ar[d, "\alpha"] \\
        \Sigma_{x:X} P(x)
    \end{tikzcd}
\end{center}

Therefore, we need something like $F(P)(F(x)) \implies P(x)$, from now on called a predicate being stable under $F$. Further if $P$ is the smallest such predicate, so for all predicates $Q$ stable under $F$ we have $P(x) \implies Q(x)$ we can then $\Sigma_{x:X} P(x)$ should be initial under the $\Sigma$ types, by just applying $F$ to the obvious function. Why this is unique needs some work still.

\begin{center}
    \begin{tikzcd}
        F(\Sigma_{x:X} P(x)) \ar[d, "\alpha"] \ar[r, "F(f)"] & F(\Sigma_{x:X} Q(x)) \ar[d, "\delta"]\\
        \Sigma_{x:X} P(x) \ar[r, "f"] & \Sigma_{x:X} Q(x)
    \end{tikzcd}
\end{center}

Finally, by building a prediacte which should be something like $\forall a:A \exists a' : A\ \ s.t. \ \ a = \tau(F(a'))$ we where $F(A) \rightarrow_{\tau} A$ is some F algebra. We should be able to build a subtype with an $F$-stable predicate that is isomorphic to $A$ and therefore get initiality for all algebras.

\newpage

\section{Recursor for Odd and Even}

In this case we can use the base recursor for $\mathbb{N}$ and build the new one from that.
We do it for odd, doing it for even should be analogous, first we define a new motive:

\begin{verbatim}
    Motive' : (n : Nat) -> (Odd(n)) x (Odd(n + 1)) -> Sort{u}
\end{verbatim}

Now assume we have a motive on odd s.t.

\begin{verbatim}
    Motive(one) 
    Motive(n) -> Motive(n + 2)
\end{verbatim}

Then for the motive above we get 

\begin{verbatim}
    Motive'(zero)
    Motive'(n) => Motive'(n + 1) 
\end{verbatim}

Now we have two cases either n is even and therefore n + 1 odd, then both Odd(n) and Odd(n+2) are trivially realized and we just pass along the Odd(n+1) function, so:

\begin{verbatim}
    Motive'(n) = Odd(n) x Odd(n + 1) -> Sort{u} 
    = Odd(n + 1) -> Sort{u} = Odd(n + 1) x Odd(n + 2) -> Sort{u} 
    = Motive'(n + 1)
\end{verbatim}

For the odd case, we use the fact that Motive(n) => Motive(n + 2)

\begin{verbatim}
    Motive'(n) = Odd(n) x Odd(n + 1) -> Sort{u} 
    = Odd(n) -> Sort{u} => Odd(n + 2) -> Sort{u} 
    = Odd(n+1) x Odd(n + 2) -> Sort{u} = Motive'(n + 1)
\end{verbatim}

What's left to show is that this covers all the Odd numbers and exactly those, which is covered by two lemmas

\begin{lemma}
    $\text{Odd}(1) \land \forall n \, (\text{Odd}(n) \implies \text{Odd}(n + 2))$
\end{lemma}

\begin{lemma}
    $\forall n \text{Odd}(n) \implies n = 1 \lor \exists n', \text{Odd}(n') \land n = n' + 2$
\end{lemma}

\begin{proof}[Sketch]\mbox{}\\
    We should be able to do this by using the fact that $+2$ is the smallest step under which Odd is closed and that Even and Odd alternate along $+1$. Then we first show that $1$ is the smallest Odd number and then recurse down the nat expression for any Odd number, showing that the number of $succ$ constructors is of the form $2n + 1$.
\end{proof}

\newpage

\section{Building Recursors}

We can make this easier by using strong induction, getting rid of the requirement to handle constructor shapes in the new motive. We then always build the motive as 
    \[NewMotive : (i : IndType) \rightarrow (p : SubType(i)) \rightarrow motive (\langle i, p \rangle).\]
Then we build the following step function:
\[ step: (i : IndType) \rightarrow ((j : IndType) \rightarrow j <_{IndOrder} i \rightarrow (pj : Subtype(j)) \rightarrow motive(\langle j, pj \rangle)) \]
\[ \rightarrow (p : Subtype(i)) \rightarrow motive \langle i , p \rangle. \]
To then be able to apply Strong induction we need two things, we first need positivity of constructors, and we secondly need "downward closure":
\begin{enumerate}
    \item $\forall \tau : Cons_{Subtype}$ we have that $\tau(\bar{x}) > \bar{x}$, basically stating that all inputs are subterms of $\tau(\bar{x})$.
    \item $\forall \langle i, p \rangle$ we have $\tau : Cons_{Subtype}$ such that $i \in im(\tau)$.
\end{enumerate}

We further need these constructors to be conservative: 
\[\forall \tau : Cons_{Subtype}, im(\tau) \subseteq SubType\]

And the motive to be "stable" along the constructors

\begin{conjecture}
    For any subtype of an inductive type which has a set of downward closed, positive and conservative constructors there exists a recursor.
\end{conjecture}

\begin{sketch}
    We build the function step by taking for any pair $\langle j , pj \rangle$ its downward element, which by positivity of the constructors is smaller than $j$.
    By the stability of the motive we then build the motive on $\langle j , pj \rangle$ from the motive on the downward element, giving us the function $step$.
    We then apply the strong induction principle of the inductive type to $NewMotive$ and therefore get a recursor on $SubType$.
\end{sketch}


\end{document}